\section{Reducing Writes with Compression}\label{sec:comppagepack}
Compression reduces DB WAF by lowering both the number of bytes written to storage and the overall dataset size.
Compressing pages before flushing to the device lowers the logical write volume.
However, if the characteristics of the underlying device are ignored, 
compression can backfire: misaligned or variable-length writes may increase read amplification or fail to reduce write and space at all.

\subsection{Page-Wise Compression}\label{sec:comp}
\module{How storage-level compression works.}
Assume a fixed database page size of 4~KiB and a buffer pool that caches at the same granularity.
Dirty pages are written to the storage device during eviction or checkpointing.
We compress each 4~KiB page \emph{independently} before persistence so it can later be read independently.
On reads, the page is decompressed before being placed in the buffer pool.

\module{Compression is difficult to apply with in-place writes.}
The SSD’s device-level write granularity complicates page-level compression for in-place systems.
\revision[R2.D4]{Compressing 4~KiB pages yields variable-length results (typically $<4$~KiB), but SSDs and filesystems commonly write in 4~KiB units.}
If we retain simple $PID \times 4\,$KiB addressing and overwrite in place, the device still writes a full 4~KiB block, eliminating any physical savings.
One could instead store compressed pages at variable offsets and maintain a PID$\rightarrow$offset mapping, but compressed sizes change as values are inserted or updated.
Offsets then drift, and relocating a page can cascade into shifting neighboring pages, making updates expensive.
Moreover, eviction follows recency rather than physical locality, so pages flushed together are rarely adjacent, hindering coalescing into aligned, sequential writes.
Consequently, many in-place DBMSs (e.g., PostgreSQL) forgo page-level compression~\cite{postgresql-toast}.
\revision[R2.D4]{Others offload compression to the filesystem (e.g., MySQL)~\cite{mysql-comp}.
However, this does not solve the fundamental issue with 4~KiB page sizes, as filesystems also issue I/Os in 4~KiB or larger units.}

\module{Optimization: \emph{easier} compression with out-of-place writes.}
With out-of-place writes, managing variable-sized compressed pages becomes far more straightforward.
For each batch of pages, we compress individual pages, write the resulting compressed batch as a large sequential I/O, and record their new offsets and compressed sizes~\cite{wiredtiger-compaction}.
This avoids the challenges of in‑place systems, making compression far easier to implement.

\module{Compression ratio determines write reduction.}
How much DB WAF drops depends on data compressibility and the algorithm.
We load several datasets into LeanStore and compare compressed vs.\ uncompressed sizes for TPC‑C (1{,}000 warehouses), YCSB‑A (100~GB)~\cite{DBLP:conf/cloud/CooperSTRS10}, and real‑world datasets. 
Using LZ4~\cite{compression/lz4} and ZSTD~\cite{compression/zstd}, compressed size (as \% of original; lower is better) range from 14\% to 49\%:

\vspace{0.2cm}

\centerline{%
\footnotesize
  \setlength\doublerulesep{0.5pt}%
  \setlength{\tabcolsep}{0.15cm}%
  \begin{tabular}{l | rrrrr}
    \toprule[1pt]\midrule[0.3pt]
    \textbf{Comp. Ratio (\%)} & \textbf{TPC-C} & \textbf{YCSB} & \textbf{Email} & \textbf{URL} & \textbf{Wiki} \\
    \midrule
    LZ4~\cite{compression/lz4} & 49.0 & 41.2 & 40.1 & 25.0 & 31.3 \\ 
    ZSTD~\cite{compression/zstd} & 36.5 & 34.4 & 27.4 & 14.5 & 17.8 \\ 
    \midrule[0.3pt]\bottomrule[1pt]
  \end{tabular}%
}
\vspace{0.2cm}


\noindent
These results confirm that most OLTP workloads will benefit from compression, as it reduces write volume and storage footprint.

\module{CPU overheads are negligible in most I/O-bound scenarios.}
Compression adds CPU work.
LZ4, for example, offers a good ratio at high speed~\cite{DBLP:journals/pacmmod/KuschewskiGNL24}, needing roughly 2--3 cycles/byte to compress and 0.5--1 cycle/byte to decompress~\cite{compression/lz4}. 
In most I/O‑bound or out‑of‑memory scenarios, the I/O savings dominate the CPU cost~\cite{DBLP:conf/sigmod/LeeM07}.
As we will show in \cref{sec:eval}, compression often improves overall performance, especially in OLTP settings.

